

\begin{tcolorbox}
\footnotesize
\texttt{<Multi-view observation><Multi-view subgoals> Task: peel vegetables. Subtask: pick up the peeler. Speed: 8000. Quality: 5. Mistake: false. Control Mode: joint.<Proprioception> }
\end{tcolorbox}
Combining all of the context information together, the example above illustrates a potential prompt that may be provided to the model.

During training, we randomly drop out each part of the prompt, which provides \MethodName{} with the flexibility to use any subset of the prompt components at test time (e.g., running with or without subgoal images).
First, we find that the model trains significantly faster when given the subgoal images --- the action prediction task essentially becomes an ``inverse dynamics'' problem inferring the robot action between the current and future frames. Thus we only add the visual subgoal images to 25\% of the examples in each batch in training. Among the examples with subgoal images we also drop out the subtask instruction $\rawtext_t$ 30\% of the time as often visual subgoal can substitute the equivalent textual subtask description in richer details. For episode metadata, we drop it entirely 15\% of the time, and additionally each component (overall speed, overall quality, and mistake label) is dropped with 5\% probability individually. We do not apply dropout for the control mode.

\section{The \texorpdfstring{\ModelSymbol{}}{pi0.7} Model and Training Recipe}
\label{sec:method}


We now discuss how we incorporate the different context in \ModelSymbol{} model by training on diverse data, as well as the details about the model architecture, training, and inference.


